\documentclass[12pt,pdflatex,authoryear]{elsarticle}

%%%%%%%%%%%%%%%%%%%%%%%%%%%%%%%%%%%%%%%%%%%%%%%%%%
%%%%%%%%%%%%%%%%%%%% PREAMBLE %%%%%%%%%%%%%%%%%%%%
%%%%%%%%%%%%%%%%%%%%%%%%%%%%%%%%%%%%%%%%%%%%%%%%%%

% word count
%TC:incbib

% references
\usepackage{natbib}

% fonts
\usepackage[default,oldstyle,scale=0.95]{opensans}
\usepackage[T1]{fontenc}
\usepackage{ae}

% colorize links
\usepackage[pdftex,hyperref,x11names]{xcolor}
\usepackage[pdftex]{hyperref}
\hypersetup{hidelinks}

% layout control
\usepackage{geometry}
\geometry{verbose,tmargin=1.25in,bmargin=1.25in,lmargin=1.1in,rmargin=1.1in}
\usepackage{fancyhdr}

% math typesetting
\usepackage{array}
\usepackage{amsmath}
\usepackage{amssymb}
\usepackage{amsfonts}
\usepackage{mathtools}
\usepackage{bm}

% float barriers at section ends
\usepackage[section]{placeins}

% tables
\usepackage{tabularx}
\usepackage{booktabs}
\usepackage{multirow}
\usepackage{seqsplit}

% caption style
\usepackage[font={small},labelfont=bf]{caption}

% footnotes at bottom
\usepackage[bottom]{footmisc}

% enumeration symbols
\usepackage{enumerate}

% graphics
\usepackage{graphicx}
\usepackage[space]{grffile}
\graphicspath{{/mnt/d/Dropbox/americium/graphics/}}

% sideways figures
\usepackage{rotating}
\usepackage{lscape}
\usepackage{pdflscape}

% spacing
\usepackage[doublespacing]{setspace}

% page template
\setlength{\headheight}{15pt}
\setlength{\headsep}{20pt}
\pagestyle{fancyplain}
\fancyhf{}
\rfoot{\fancyplain{}{\thepage}}

% customizations
\newcommand{\mbs}[1]{\boldsymbol{#1}}
\newcommand{\pkg}[1]{{\fontseries{b}\selectfont #1}}

%%%%%%%%%%%%%%%%%%%%%%%%%%%%%%%%%%%%%%%%%%%%%%%%%%
%%%%%%%%%%%%%%%%%%%% DOCUMENT %%%%%%%%%%%%%%%%%%%%
%%%%%%%%%%%%%%%%%%%%%%%%%%%%%%%%%%%%%%%%%%%%%%%%%%

% remove elsevier preprint note
\makeatletter
\def\ps@pprintTitle{%
 \let\@oddhead\@empty
 \let\@evenhead\@empty
 \def\@oddfoot{}%
 \let\@evenfoot\@oddfoot}
\makeatother


\begin{document}

\thispagestyle{empty}
\begin{frontmatter}

\title{Supplementary Information for\\
``When America Stops Bridging: Diplomatic Discontinuity in UN Voting''}

\author[strath]{Max Gallop}
\ead{max.gallop@strath.ac.uk}
\author[msu]{Shahryar Minhas}
\ead{minhassh@msu.edu}

\address[strath]{Department of Government and Public Policy, University of Strathclyde}
\address[msu]{Department of Political Science, Michigan State University}

\end{frontmatter}

\newpage\setcounter{page}{1}

\noindent This supplementary information gives the full specification of the
snap-shift latent-space model, the Gibbs sampler, the identification strategy,
the simulation study validating recovery, robustness checks, and the correlates
of the 2025 snap referenced in the main text.

\newpage
\appendix
\section{The Snap-Shift Latent-Space Model}
\label{app:model}

This appendix gives the complete specification of the model, the prior
distributions, the Gibbs sampler, the treatment of identification, and the
simulation study that validates recovery. We write the model for a single
relation; the three relations are fit separately.

\subsection{Data and notation}

For a relation observed over years $t=1,\dots,T$ on a common actor set
$\{1,\dots,n\}$, let $Y_t$ denote the $n\times n$ sociomatrix in year $t$, with
$y_{t,ij}$ the directed tie from $i$ to $j$ and $y_{t,ii}$ undefined. Actor
composition may vary across years; an actor absent in year $t$ contributes no
observations that year and its effects for that year are drawn from their priors.
The latent rank $R$ is fixed in advance ($R=2$ in the main analysis). $I_R$ is
the $R\times R$ identity.

\subsection{Focal actor set}

The focal rupture index and focal-actor summaries use a theory-selected set of
order-relevant actors: the hegemon, core allies, strategic rivals, and pivotal
hedgers. Table~\ref{tab:actor-set} lists the exact denominator. States outside
this set remain in the fitted networks; they are excluded only from the focal
summary index so that the index answers the theoretical question about
weight-bearing actors rather than averaging over all observed states.

The category rule is substantive rather than outcome-driven. The hegemon is the
United States. Core allies are long-standing U.S.-aligned democratic partners in
the Atlantic and Indo-Pacific wings whose diplomatic position bears directly on
whether the allied core follows, fragments, or holds together after U.S.
defection. Strategic rivals are states that persistently occupy adversarial or
counter-order positions relative to the U.S.-led order in the post--Cold War
period. Pivotal hedgers are regionally or economically important states with
elastic cross-bloc strategies, so their movement indicates how far a rupture
propagates beyond the hegemon-allied-rival core. The denominator is therefore
not a claim that other states are irrelevant; it is the subset on which the
paper's theory makes the sharpest configurational predictions.

\begin{table}[ht]
\centering
\caption{Theory-selected focal actor categories used in the rupture-index and
focal-actor summaries.}
\label{tab:actor-set}
\small
\begin{tabularx}{0.95\textwidth}{@{}l X@{}}
\toprule
\textbf{Category} & \textbf{Actors} \\
\midrule
Hegemon & USA \\
Core allies & FRA, DEU, GBR, ITA, CAN, JPN, AUS, NLD, ESP, KOR, BEL, NOR, DNK, POL \\
Strategic rivals & RUS, CHN, PRK, IRN, BLR, SYR, VEN, CUB, NIC \\
Pivotal hedgers & IND, TUR, BRA, ZAF, IDN, SAU, MEX, EGY, ARE, VNM, THA \\
\bottomrule
\end{tabularx}
\end{table}

Table~\ref{tab:denominator-contrasts} repeats the index for wider and
lower-exposure denominators. The focal actors have the highest 2025
concentration, but the all-state and non-focal contrasts are also elevated. The
implication is not that peripheral actors are inert; it is that the 2025 focal
index is not produced by choosing a denominator that excludes a volatile
periphery or by counting only the United States.

\begin{table}[ht]
\centering
\caption{Denominator contrasts for the UN-voting rupture index. The index is the
mean posterior snap probability for the listed actor set. The low-exposure
contrast is the non-focal bottom quartile of 2024 U.S. trade dependence among
states with complete covariates.}
\label{tab:denominator-contrasts}
\small
\begin{tabular}{@{}l r r r r@{}}
\toprule
\textbf{Denominator} & \textbf{N in index} & \textbf{2024 index} &
\textbf{2025 index} & \textbf{2025 rank} \\
\midrule
Theory-selected focal actors & 35 & 0.030 & 0.956 & 1 \\
All observed states & 191 & 0.146 & 0.776 & 1 \\
Observed non-focal states & 157 & 0.173 & 0.732 & 1 \\
Low-exposure non-focal contrast & 34 & 0.228 & 0.742 & 1 \\
\bottomrule
\end{tabular}
\end{table}

The focal denominator contains $35$ theory-selected actors. Of these, $34$ are
observed in the 2025 UN-voting layer and all $34$ have posterior 2025 snap
probability above $0.5$. Venezuela is the focal actor not observed in the 2025
UN-voting layer; it is retained in the denominator with posterior snap
probability $0.068$, making the reported focal index conservative. Excluding
that unobserved focal actor gives a 2025 focal index of $0.982$.

Table~\ref{tab:leaveout-index} reports the leave-out checks referenced in the
main text. These are aggregation checks using the main posterior snap
probabilities; they are not separate no-U.S. refits. Their purpose is narrower:
to show that the focal rupture index remains near one when the headline actor,
the two largest rivals, or all strategic rivals are removed from the denominator.

\begin{table}[ht]
\centering
\caption{Leave-out focal rupture indices in the main UN-voting fit. Entries recompute the reported index over different actor sets using the same posterior snap probabilities; they are aggregation checks, not separate no-U.S. refits.}
\label{tab:leaveout-index}
\small
\begin{tabularx}{\textwidth}{@{}X r r r@{}}
\toprule
\textbf{Index denominator} & \textbf{2024} & \textbf{2025} & \textbf{Change} \\
\midrule
Full focal index & 0.030 & 0.956 & +0.926 \\
Excluding United States & 0.031 & 0.955 & +0.923 \\
Excluding United States, Russia, and China & 0.033 & 0.952 & +0.919 \\
Excluding United States and all strategic rivals & 0.006 & 0.975 & +0.969 \\
\bottomrule
\end{tabularx}
\end{table}


\subsection{Network construction checks}

All three modeled relations are aligned to the same UN member-state actor grid
and the common 1989--2025 year span before estimation. Self-ties are undefined
in the model and are coded as missing in every relation. Table~\ref{tab:data-qc}
summarizes the aligned arrays used by the sampler. The sanctions layer is the
state-designator panel after this alignment; supranational designators such as
EU and UN are present in the source parsing but are not actors in the fitted
state-by-state network.

UN voting comes from the United Nations Digital Library bulk voting file, which
includes the February 24, 2025 Ukraine resolution used in the article's
motivating example. Trade uses the IMF International Merchandise Trade
Statistics API annual FOB exports series, transformed into bilateral
merchandise-trade networks. Sanctions use the OpenSanctions consolidated
sanctions-targets file, parsed into dated state-designator to target-country
records. Because this is an event-count layer, non-self dyads with no observed
designation in a year are coded as zero rather than missing. The trade and
sanctions layers therefore use the annual evidence available at the time of
analysis; they are not complete post-2025 propagation data.

\begin{table}[ht]
\centering
\caption{Construction checks for aligned modeled arrays. Observed off-diagonal
counts summarize non-missing dyads by year; fully unobserved actors are actors
with no observed row or column entries in a given relation-year.}
\label{tab:data-qc}
\scriptsize
\begin{tabular}{@{}l l r r r r r@{}}
\toprule
\textbf{Relation} & \textbf{Years} & \textbf{Actors} &
\textbf{Min. dyads} & \textbf{Median dyads} &
\textbf{Max. unobs. actors} &
\textbf{Non-missing diag.} \\
\midrule
UN voting & 1989--2025 & 200 & 23562 & 36290 & 46 & 0 \\
Trade & 1989--2025 & 200 & 12478 & 24434 & 43 & 0 \\
Sanctions & 1989--2025 & 200 & 39800 & 39800 & 0 & 0 \\
\bottomrule
\end{tabular}
\end{table}

\subsection{Likelihood and data augmentation}

Each observed tie is linked to a latent Gaussian variable $z_{t,ij}$ through a
family-specific map, and the additive and multiplicative structure is placed on
$z$:
\begin{equation}
  z_{t,ij} \;=\; \mu_t + a_{t,i} + b_{t,j}
    + \mbs{u}_{t,i}^\top \mbs{v}_{t,j} + \epsilon_{t,ij},
  \qquad
  \begin{pmatrix}\epsilon_{t,ij}\\ \epsilon_{t,ji}\end{pmatrix}
  \sim \mathcal{N}\!\left(\mbs{0},\,
  \begin{pmatrix} 1 & \rho \\ \rho & 1 \end{pmatrix}\right).
  \label{eq:srm}
\end{equation}
The within-dyad correlation $\rho$ encodes reciprocity. The map from $z_{t,ij}$ to
$y_{t,ij}$ depends on the family: for the Gaussian case (importance-weighted UN
agreement, and log trade) $y_{t,ij}=z_{t,ij}$ up to a scale; for the Poisson case
(sanctions counts) $y_{t,ij}\mid z_{t,ij}\sim\mathrm{Poisson}(e^{z_{t,ij}})$.
In each non-Gaussian case a data-augmentation step draws $z_{t,ij}$ from its full
conditional given $y_{t,ij}$, after which every remaining update is conjugate
Gaussian. For symmetric relations (UN agreement, total trade) we impose
$\mbs{v}_{t,i}=\mbs{u}_{t,i}$ and $b_{t,i}=a_{t,i}$, so a single set of latent
positions describes each actor.

\subsection{Additive effects}

The additive effects follow the social-relations model. Stacking
$(a_{t,i},b_{t,i})$, we take
\begin{equation}
  (a_{t,i}, b_{t,i})^\top \sim \mathcal{N}(\mbs{0}, \Sigma_{ab}),
  \qquad
  \Sigma_{ab} \sim \mathrm{Inverse\text{-}Wishart}(\eta_0, \eta_0 \Sigma_0),
\end{equation}
which lets a state's propensity to send and to receive ties covary. The baseline
$\mu_t$ absorbs the overall density of $Y_t$.

\subsection{Latent positions and the snap-shift transition}

The multiplicative factors carry the dynamics. For $t=1$, each
$\mbs{u}_{1,i}\sim\mathcal{N}(\mbs{0}, \sigma_{u,0}^2 I_R)$; no snap indicator is
defined for the first year because no previous position exists from which the
actor could drift. For $t\ge 2$, the transition is the two-component mixture of
the main text:
\begin{equation}
  \mbs{u}_{t,i} \mid \delta_{t,i}, \mbs{u}_{t-1,i}
  \;\sim\;
  (1-\delta_{t,i})\,
    \mathcal{N}\!\left(\rho_{uv}\,\mbs{u}_{t-1,i}, \sigma_u^2 I_R\right)
  \;+\;
  \delta_{t,i}\,
    \mathcal{N}\!\left(\mbs{0}, \kappa^2 I_R\right),
  \label{eq:transition}
\end{equation}
with $\delta_{t,i}\sim\mathrm{Bernoulli}(\pi)$ independently across eligible $i$
and $t$, and $\kappa^2\gg\sigma_u^2$. The drift component is the smooth AR(1)
transition used by standard dynamic latent-position models; the snap component is
a diffuse re-draw that discards the past. The persistence
$\rho_{uv}$, innovation variance $\sigma_u^2$, jump scale $\kappa$, and rate $\pi$ are
common across states, so that ``how much a typical state drifts'' and ``how large
and how common a jump is'' are estimated from the whole panel rather than imposed.

\subsection{Priors}

We complete the model with weakly informative priors:
\begin{align*}
  \rho_{uv} &\sim \mathcal{N}(m_\rho, s_\rho^2)\ \text{truncated to } (-1,1),
    & \sigma_u^{-2} &\sim \mathrm{Gamma}(\nu_u/2, \nu_u s_u^2/2), \\
  \pi &\sim \mathrm{Beta}(a_\pi, b_\pi),\ b_\pi \gg a_\pi,
    & \rho &\sim \mathrm{Uniform}(-1,1).
\end{align*}
The jump scale $\kappa$ is fixed on the standardized latent scale; the main
analysis uses $\kappa=2$, making the snap prior diffuse relative to ordinary
drift innovations estimated by $\sigma_u^2$. The $\mathrm{Beta}$ prior on $\pi$
with $b_\pi\gg a_\pi$ is the formal statement that snaps are rare events. The
main analysis uses $a_\pi=1, b_\pi=999$, a prior mean of $0.001$; this conservative
calibration is chosen because drift-only simulations under weaker rare-event
priors generated too many false snaps.

\subsection{Metropolis-within-Gibbs sampler}

One sweep of the sampler cycles through the following conditional updates.
\begin{enumerate}
\item \textbf{Augmentation.} For non-Gaussian families, update $z_{t,ij}$ given
$y_{t,ij}$ and the current linear predictor. The Gaussian and binary cases use
their standard truncated-Gaussian updates; the Poisson sanctions model uses a
Metropolis--Hastings update for the augmented $z$ values.
\item \textbf{Baseline and additive effects.} Draw $\mu_t$, then
$(a_{t,i}, b_{t,i})$ from their Gaussian full conditionals given the SRM
covariance, and update $\Sigma_{ab}$ from its inverse-Wishart conditional.
\item \textbf{Latent positions and indicators.} For each $t\ge 2$ and each $i$,
form the residual $r_{t,ij} = z_{t,ij} - \mu_t - a_{t,i} - b_{t,j}$, and the
Gaussian information $V_{t}^\top V_{t}$ and $V_{t}^\top r_{t,i\cdot}$ that the
year-$t$ ties carry about $\mbs{u}_{t,i}$. Under each regime the position has a
Gaussian conditional update, a precision-weighted combination of this information
with the regime's prior. The drift prior is centered at
	$\rho_{uv}\,\mbs{u}_{t-1,i}$ with precision $\sigma_u^{-2}$; the snap prior is
centered at $\mbs{0}$ with precision $\kappa^{-2}$. A forward contribution from
year $t+1$ is included only when the current value of $\delta_{t+1,i}$ places the
next transition in the drift regime; if year $t+1$ is itself a snap, it is
conditionally independent of $\mbs{u}_{t,i}$. Each regime implies a
closed-form Gaussian log marginal likelihood $\log\mathcal{L}_{\text{drift}}$ and
$\log\mathcal{L}_{\text{snap}}$ for $r_{t,i\cdot}$. Draw
\[
  \delta_{t,i}\sim\mathrm{Bernoulli}\!\Big(
  \sigma\big(\operatorname{logit}\pi
  + \log\mathcal{L}_{\text{snap}}-\log\mathcal{L}_{\text{drift}}\big)\Big),
\]
clamping the log-odds to $[-20,20]$ for numerical stability, and then draw
$\mbs{u}_{t,i}$ from the full conditional of the selected regime.
\item \textbf{Snap rate.} With $S=\sum_{t\ge 2,i}\delta_{t,i}$ snaps among the
$M$ eligible actor-years, draw $\pi\sim\mathrm{Beta}(a_\pi+S,\, b_\pi+M-S)$.
\item \textbf{Hyperparameters.} Update $\rho$ (reciprocity) by Metropolis--Hastings;
update $\rho_{uv}$ (persistence) and $\sigma_u^2$ (innovation variance) from their
conditionals, computed over the \emph{drift} transitions only, so that snapped
increments do not inflate the estimate of how fast states ordinarily drift.
\end{enumerate}
These restrictions are required for the snap probability to have the rupture
interpretation used in the main text. If snapped increments were included in the
drift hyperparameter updates, the ordinary drift variance would learn from the
very breaks it is meant to distinguish. If the forward AR term were retained when
$\delta_{t+1,i}=1$, a reset in year $t+1$ would still force year $t$ to predict it.
The implemented sampler therefore makes a snap a conditional reset rather than an
oversized AR innovation.
For the final observed year, no forward AR contribution exists by construction.
The 2025 snap probabilities are therefore boundary-year estimates based on how
well the 2025 relational profile is explained by continuity from 2024 versus a
diffuse reset. They are not two-sided smoothed classifications that use future
information. The simulation below includes a terminal planted snap and a
drift-only terminal-year check for this reason.
For UN voting and trade, we run the sampler for $6000$ iterations after $3000$
burn-in, thinning to every tenth draw. Because the directed Poisson sanctions
fit mixes more slowly, we run that fit for $30000$ iterations after $15000$
burn-in, with the same thinning interval. Convergence is assessed by trace plots
and effective sample sizes of the saved scalar chains, especially $\pi$,
$\rho_{uv}$, and $\sigma_u^2$. The
snap probabilities are posterior means of retained indicator draws. For every
focal rupture index, we also average the retained indicators within each draw and
summarize the resulting draw-level index. The reported intervals and rank
probabilities therefore reflect posterior uncertainty in the snap assignments
rather than treating the posterior mean probability matrix as fixed.
Table~\ref{tab:snap-index-uncertainty} reports these draw-level summaries for the
three main relation-level comparisons. The UN-voting fit is the primary
empirical fit in the article; the trade fit is a cross-layer informative null;
and the sanctions fit is a directed coercive-statecraft diagnostic. The directed
Poisson fit mixes more slowly on the saved scalar chains, so the sanctions layer
does not carry the main claim. The reported sanctions estimates show whether
coercive relations move on the same schedule and side of the network as UN voting
and trade.

\begin{table}[ht]
\centering
\caption{Draw-level uncertainty for the 2025 focal rupture indices. Entries average retained snap-indicator draws within the theory-selected focal actor set before summarizing the index. Rank probabilities compare 2025 with all eligible transition years in the same relation and index side.}
\label{tab:snap-index-uncertainty}
\small
\begin{tabularx}{\textwidth}{@{}l l r r r r@{}}
\toprule
\textbf{Relation} & \textbf{Index side} & \textbf{Mean} & \textbf{95\% interval} & \textbf{Pr(rank 1)} & \textbf{Pr(rank $\leq$ 3)} \\
\midrule
UN voting & State/sender & 0.956 & [0.914, 1.000] & 1.00 & 1.00 \\
Trade & State/sender & 0.023 & [0.000, 0.086] & 0.17 & 0.39 \\
Sanctions & State/sender & 0.089 & [0.057, 0.143] & 0.35 & 0.85 \\
Sanctions & Receiver/target & 0.379 & [0.286, 0.486] & 1.00 & 1.00 \\
Sanctions & Sender or receiver & 0.433 & [0.343, 0.543] & 1.00 & 1.00 \\
\bottomrule
\end{tabularx}
\end{table}


Table~\ref{tab:mcmc-diagnostics} reports the scalar diagnostics for the three
main fits. The saved result objects also record the \pkg{lame} version and MD5
hashes for the local snap-shift source files; the reproduction scripts check
those hashes before refreshing manuscript numbers. A public replication archive
should vendor or otherwise pin the same package source rather than rely on a
mutable local branch. Because the local \pkg{lame} branch can receive later
comment or documentation cleanup, Table~\ref{tab:mcmc-diagnostics}
distinguishes recorded source hashes from the current local-source match. The
recorded hash is the relevant provenance for the fitted objects; any public
replication release should freeze the exact package source used for the final
fits.

\begin{table}[ht]
\centering
\caption{Scalar MCMC diagnostics for the three main snap-shift fits. The scan column reports total scans, burn-in, and thinning interval; retained draws are posterior draws after burn-in and thinning. Parentheses report posterior standard deviations. ESS is the minimum effective sample size across the saved scalar chains for $\pi$, $\rho_{uv}$, and $\sigma_u$. Recorded source indicates that the saved fit contains the snap-shift source hash; local match reports whether the currently mutable local \pkg{lame} branch matches that saved hash.}
\label{tab:mcmc-diagnostics}
\scriptsize
\resizebox{\textwidth}{!}{%
\begin{tabular}{@{}l l l r r r r r l l@{}}
\toprule
\textbf{Relation} & \textbf{Family} & \textbf{Scans/burn/thin} & \textbf{Draws} & $\bar{\pi}$ & $\rho_{uv}$ & $\sigma_u$ & \textbf{Min. ESS} & \textbf{Recorded source} & \textbf{Local match} \\
\midrule
UN voting & normal &  6000/ 3000/10 &  600 & 0.369 (0.007) & 0.986 (0.001) & 0.011 (0.000) & 154 & yes & no \\
Trade & normal &  6000/ 3000/10 &  600 & 0.012 (0.004) & 0.990 (0.000) & 0.008 (0.000) & 119 & yes & no \\
Sanctions & poisson & 30000/15000/10 & 3000 & 0.024 (0.003) & 0.000 (0.000) & 0.007 (0.000) & 32 & yes & yes \\
\bottomrule
\end{tabular}
}
\end{table}


\subsection{Identification and rotation invariance}

The multiplicative term is invariant to a simultaneous orthogonal rotation of all
latent positions, $\mbs{u}_{t,i}\mapsto Q\mbs{u}_{t,i}$, $\mbs{v}_{t,j}\mapsto
Q\mbs{v}_{t,j}$ for any orthogonal $Q$. The AR(1) prior couples adjacent years and
so partially pins the rotation across time, but a snap, by re-drawing from an
isotropic prior, relaxes exactly this coupling. We therefore (i) draw the snap
indicators from the rotation-invariant marginal likelihoods above, which depend on
the positions only through inner products, and (ii) report all substantive spatial
claims, including bloc cohesion and the distance between the hegemon and its allies, in terms
of rotation- and translation-invariant latent distances. The cooperation-space
maps in the main text use posterior mean latent positions and apply a common-actor
Procrustes rotation and translation to align 2025 to 2024. Posterior uncertainty
in the positions is summarized by the per-coordinate posterior standard
deviation, accumulated alongside the posterior mean. We do not interpret the
raw coordinate axes substantively.

\subsection{What the snap probabilities add}

The model produces an actor-year posterior surface as well as a system-level
change score. Figure~\ref{fig:snap-history} reports that surface for the
theory-selected focal actors over the full post-Cold War series. The heatmap
shows the limit of the focal rupture index: some earlier years contain high snap
probabilities for particular states or categories, while 2025 is distinctive
because high probabilities arrive broadly across the hegemon, allies, rivals, and
hedgers. The model is not assigning a label to the final year of the series; it
estimates where discontinuity is and is not present for each actor-year.

\begin{figure}[ht]
\centering
\includegraphics[width=\textwidth]{snap_history_heatmap.png}
\caption{Posterior snap probabilities for theory-selected focal actors in the
UN-voting fit, 1990--2025. Rows are ordered by actor category; columns are
eligible transition years. Darker red indicates higher posterior probability
that the actor-year is better described as a discontinuous reset than as smooth
drift. The saturated 2025 column means that the broad focal set, not only the
United States, is classified as snapping in 2025. Dashed vertical lines mark
benchmark moments used in the historical calibration.}
\label{fig:snap-history}
\end{figure}

Tables~\ref{tab:episode-typology} and~\ref{tab:country-snap-history} summarize the
same surface in tabular form. Table~\ref{tab:episode-typology} shows that earlier
episodes are real but partial: some are rival- and hedger-heavy, some are
allied-heavy, and some are broad but moderate. Table~\ref{tab:country-snap-history}
shows that many focal actors have high-probability snap years before 2025. The
distinctive feature of 2025 is therefore synchronization across actor histories,
not the mere existence of a large U.S. movement.

Table~\ref{tab:raw-un-diagnostics} reports descriptive UN-voting diagnostics for
2024 and 2025. These summaries do not replace the model: the annual UN agenda
changes, so the table is not a same-resolution change score. It serves as a
construct-validity check. The raw record shows the U.S. no-share rising sharply,
U.S. agreement falling with allies, rivals, and hedgers, and allied internal
agreement remaining high without U.S. ties.

Table~\ref{tab:un-agenda-diagnostics} reports the corresponding agenda
diagnostic. The 2025 roll-call agenda is unusually large in the cached data:
there are 192 resolutions in 2025 and 95 in 2024, and the 2025 resolutions are
less competitive on average. The issue mix also changes. This is not a nuisance
to be waved away. It means the annual UN-voting result is evidence about the
observed 2025 diplomatic agenda, not a matched-resolution estimate of how the
same issues would have been voted on under continued U.S. role performance. The
construction checks below show that the result is not an artifact of
divisiveness weighting or the most polarizing resolutions, but a fully
issue-matched design remains a natural next robustness check.

Table~\ref{tab:un-issue-agreement} adds a reduced-form issue-stratified
diagnostic. Within every broad agenda group, raw U.S.-ally agreement falls from
2024 to 2025. Allied internal agreement, by contrast, is flat or higher in every
group. This does not replace an issue-matched snap-model refit, because each
group still contains different resolutions across years. It does show that the
basic descriptive pattern is not confined to one agenda class: U.S.-ally
separation and allied cohesion appear within Ukraine/Russia votes,
Palestine/Israel votes, human-rights votes, development and economic votes,
security and disarmament votes, and the residual institutional category.

\begin{table}[ht]
\centering
\caption{Selected diplomatic snap episodes in the UN-voting layer. Category entries are posterior mean snap probabilities for the focal actors in that category.}
\label{tab:episode-typology}
\scriptsize
\begin{tabularx}{\textwidth}{@{}r r r r r r r X@{}}
\toprule
\textbf{Year} & \textbf{Index} & \textbf{U.S.} & \textbf{Allies} & \textbf{Rivals} & \textbf{Hedgers} & \textbf{$>0.5$} & \textbf{Signature} \\
\midrule
1991 & 0.406 & 0.110 & 0.431 & 0.408 & 0.400 & 12 & post-Cold War adjustment; allies, rivals, and hedgers move more than the United States \\
1998 & 0.511 & 0.627 & 0.399 & 0.681 & 0.503 & 17 & late-1990s rival and hedger re-coordination, with Turkey, China, Belarus, and Japan prominent \\
2001 & 0.465 & 0.457 & 0.394 & 0.502 & 0.525 & 10 & post-9/11 and Bush-era re-coordination; rivals and hedgers are most exposed \\
2009 & 0.619 & 0.843 & 0.678 & 0.531 & 0.595 & 24 & financial-crisis-era diplomatic re-coordination; broad but less synchronized than 2025 \\
2013 & 0.593 & 0.630 & 0.737 & 0.474 & 0.504 & 22 & pre-Crimea and Arab Spring-era bloc repositioning; allies and rivals are high \\
2019 & 0.498 & 0.505 & 0.526 & 0.512 & 0.451 & 18 & Trump I-era re-coordination; broad but moderate relative to 2025 \\
2023 & 0.544 & 0.740 & 0.585 & 0.446 & 0.554 & 19 & Russia-invasion aftermath; United States, allies, and hedgers are high \\
2025 & 0.956 & 1.000 & 1.000 & 0.896 & 0.944 & 34 & bloc-preserving hegemonic exit; all focal categories snap together \\
\bottomrule
\end{tabularx}
\end{table}


\begin{sidewaystable}[ht]
\centering
\caption{Country-level snap histories for theory-selected focal actors in the UN-voting layer. The final column lists the four highest posterior snap-probability years for each actor.}
\label{tab:country-snap-history}
\scriptsize
\begin{tabularx}{\textwidth}{@{}l l r r r X@{}}
\toprule
\textbf{Actor} & \textbf{Category} & \textbf{Max. year} & \textbf{Max. prob.} & \textbf{2025 prob.} & \textbf{Top snap years} \\
\midrule
USA & Hegemon & 2025 & 1.000 & 1.000 & 2025 (1.00); 1996 (0.85); 2009 (0.84); 2023 (0.74) \\
FRA & Core Allies & 2025 & 1.000 & 1.000 & 2025 (1.00); 2013 (0.85); 1992 (0.84); 2021 (0.75) \\
DEU & Core Allies & 2025 & 1.000 & 1.000 & 2025 (1.00); 2013 (0.89); 2009 (0.81); 2007 (0.79) \\
GBR & Core Allies & 2025 & 1.000 & 1.000 & 2025 (1.00); 1996 (0.89); 2013 (0.81); 2021 (0.79) \\
ITA & Core Allies & 2025 & 1.000 & 1.000 & 2025 (1.00); 2021 (0.77); 2009 (0.77); 2013 (0.66) \\
CAN & Core Allies & 2007 & 1.000 & 1.000 & 2007 (1.00); 2011 (1.00); 2025 (1.00); 2009 (0.99) \\
JPN & Core Allies & 1993 & 1.000 & 1.000 & 1993 (1.00); 2025 (1.00); 1992 (0.99); 2011 (0.93) \\
AUS & Core Allies & 1991 & 1.000 & 1.000 & 1991 (1.00); 2004 (1.00); 2025 (1.00); 2013 (0.81) \\
NLD & Core Allies & 2025 & 1.000 & 1.000 & 2025 (1.00); 2007 (0.77); 2009 (0.73); 2013 (0.64) \\
ESP & Core Allies & 2025 & 1.000 & 1.000 & 2025 (1.00); 1994 (0.88); 2011 (0.85); 2023 (0.80) \\
KOR & Core Allies & 1991 & 1.000 & 1.000 & 1991 (1.00); 2002 (1.00); 2025 (1.00); 2016 (0.98) \\
BEL & Core Allies & 2025 & 1.000 & 1.000 & 2025 (1.00); 2023 (0.78); 2013 (0.76); 2009 (0.70) \\
NOR & Core Allies & 2025 & 1.000 & 1.000 & 2025 (1.00); 1991 (0.96); 2021 (0.84); 2013 (0.65) \\
DNK & Core Allies & 2025 & 1.000 & 1.000 & 2025 (1.00); 2013 (0.79); 2007 (0.75); 2009 (0.70) \\
POL & Core Allies & 1990 & 1.000 & 1.000 & 1990 (1.00); 2025 (1.00); 1991 (0.97); 2007 (0.78) \\
RUS & Strategic Rivals & 1992 & 1.000 & 1.000 & 1992 (1.00); 2000 (1.00); 2004 (1.00); 2006 (1.00) \\
CHN & Strategic Rivals & 2022 & 1.000 & 1.000 & 2022 (1.00); 2025 (1.00); 1998 (0.99); 2000 (0.98) \\
PRK & Strategic Rivals & 1991 & 1.000 & 1.000 & 1991 (1.00); 2025 (1.00); 2017 (0.98); 2022 (0.98) \\
IRN & Strategic Rivals & 2025 & 1.000 & 1.000 & 2025 (1.00); 2022 (0.98); 1996 (0.68); 1998 (0.68) \\
BLR & Strategic Rivals & 1991 & 1.000 & 1.000 & 1991 (1.00); 2000 (1.00); 2025 (1.00); 1998 (0.99) \\
SYR & Strategic Rivals & 2025 & 1.000 & 1.000 & 2025 (1.00); 2022 (0.93); 1996 (0.93); 1998 (0.81) \\
VEN & Strategic Rivals & 1995 & 0.883 & 0.068 & 1995 (0.88); 2009 (0.80); 2022 (0.68); 2013 (0.54) \\
CUB & Strategic Rivals & 2025 & 1.000 & 1.000 & 2025 (1.00); 1996 (0.83); 2022 (0.63); 1998 (0.62) \\
NIC & Strategic Rivals & 1993 & 1.000 & 1.000 & 1993 (1.00); 2001 (1.00); 2002 (1.00); 2007 (1.00) \\
IND & Pivotal Hedgers & 2025 & 1.000 & 1.000 & 2025 (1.00); 2005 (0.99); 2006 (0.99); 1994 (0.98) \\
TUR & Pivotal Hedgers & 1998 & 1.000 & 1.000 & 1998 (1.00); 2025 (1.00); 1994 (0.99); 2021 (0.98) \\
BRA & Pivotal Hedgers & 2009 & 1.000 & 0.825 & 2009 (1.00); 2019 (1.00); 2023 (1.00); 2018 (1.00) \\
ZAF & Pivotal Hedgers & 1999 & 1.000 & 1.000 & 1999 (1.00); 2001 (1.00); 2002 (1.00); 2025 (1.00) \\
IDN & Pivotal Hedgers & 2025 & 1.000 & 1.000 & 2025 (1.00); 1993 (0.81); 2009 (0.62); 1991 (0.57) \\
SAU & Pivotal Hedgers & 2025 & 1.000 & 1.000 & 2025 (1.00); 2001 (0.73); 2023 (0.65); 1998 (0.51) \\
MEX & Pivotal Hedgers & 2010 & 1.000 & 0.743 & 2010 (1.00); 2011 (1.00); 2005 (0.99); 2006 (0.98) \\
EGY & Pivotal Hedgers & 2025 & 1.000 & 1.000 & 2025 (1.00); 1999 (0.69); 2009 (0.66); 1995 (0.56) \\
ARE & Pivotal Hedgers & 2001 & 0.998 & 0.998 & 2001 (1.00); 2025 (1.00); 2000 (0.81); 2009 (0.78) \\
VNM & Pivotal Hedgers & 2025 & 1.000 & 1.000 & 2025 (1.00); 1993 (0.94); 2022 (0.69); 1995 (0.67) \\
THA & Pivotal Hedgers & 2009 & 0.988 & 0.818 & 2009 (0.99); 2025 (0.82); 2023 (0.72); 2001 (0.68) \\
\bottomrule
\end{tabularx}
\end{sidewaystable}


\begin{table}[ht]
\centering
\caption{Raw UN-voting diagnostics around 2025. Entries are descriptive annual-agenda summaries, not model estimates and not same-resolution before--after comparisons. Agreement is Lijphart three-category agreement on yes, no, and abstain votes.}
\label{tab:raw-un-diagnostics}
\scriptsize
\begin{tabularx}{\textwidth}{@{}X r r r X@{}}
\toprule
\textbf{Raw diagnostic} & \textbf{2024} & \textbf{2025} & \textbf{Change} & \textbf{Interpretation} \\
\midrule
U.S. share voting yes & 0.340 & 0.052 & -0.288 & U.S. vote profile in the annual UN agenda \\
U.S. share voting no & 0.543 & 0.906 & +0.364 & U.S. vote profile in the annual UN agenda \\
U.S. share abstaining & 0.117 & 0.042 & -0.075 & U.S. vote profile in the annual UN agenda \\
U.S. agreement with core allies & 0.709 & 0.238 & -0.472 & direct U.S.-ally divergence \\
U.S. agreement with strategic rivals & 0.249 & 0.112 & -0.137 & U.S. does not simply join the rival voting profile \\
U.S. agreement with pivotal hedgers & 0.384 & 0.094 & -0.290 & U.S. diverges from pivotal hedgers as well \\
Agreement within core allies & 0.925 & 0.953 & +0.027 & allied cohesion excluding U.S. ties \\
Europe--Indo-Pacific allied agreement & 0.915 & 0.942 & +0.028 & Europe and Indo-Pacific allies remain close without U.S. ties \\
Core-allied agreement with strategic rivals & 0.507 & 0.759 & +0.252 & cross-bloc agreement also changes on the annual agenda \\
Core-allied agreement with pivotal hedgers & 0.641 & 0.816 & +0.175 & allied relation to hedgers changes on the annual agenda \\
\bottomrule
\end{tabularx}
\end{table}


\begin{table}[!htbp]
\centering
\begingroup
\scriptsize
\caption{Descriptive UN-voting agenda diagnostics, 2024--2025. Agenda groups are rule-based summaries from resolution titles, agenda titles, and UN subject tags. The categories are diagnostic rather than mutually exhaustive substantive codings; the model uses the full annual voting network.}
\label{tab:un-agenda-diagnostics}
\begin{tabularx}{\textwidth}{@{}rrrrrrrrrr@{}}
\toprule
Year & Res. & Turnout & Compet. & Ukr./Rus. & Pal./Isr. & Rights & Dev./econ. & Sec./disarm. & Other \\
\midrule
2024 & 95 & 0.920 & 0.193 & 2 & 17 & 13 & 10 & 33 & 20 \\
2025 & 192 & 0.901 & 0.110 & 5 & 16 & 29 & 57 & 37 & 48 \\
\bottomrule
\end{tabularx}
\endgroup
\end{table}


\begin{table}[!htbp]
\centering
\begingroup
\scriptsize
\caption{Issue-stratified raw UN-voting agreement diagnostics, 2024--2025. Entries are Lijphart three-category agreement averaged within broad agenda groups. The diagnostic asks whether U.S.-ally agreement falls and allied internal agreement holds inside issue strata; it is not a snap-model refit.}
\label{tab:un-issue-agreement}
\begin{tabularx}{\textwidth}{@{}>{\raggedright\arraybackslash}Xrrrrrrr@{}}
\toprule
Agenda group & Res. 24/25 & U.S.-ally 24 & U.S.-ally 25 & Change & Ally internal 24 & Ally internal 25 & Change \\
\midrule
Ukraine/Russia & 2/5 & 1.000 & 0.500 & -0.500 & 1.000 & 1.000 & 0.000 \\
Palestine/Israel & 17/16 & 0.248 & 0.080 & -0.168 & 0.888 & 0.942 & 0.053 \\
Human rights/self-determination & 13/29 & 0.890 & 0.385 & -0.505 & 0.948 & 0.968 & 0.021 \\
Development/economic & 10/57 & 0.511 & 0.094 & -0.417 & 0.960 & 0.970 & 0.010 \\
Security/disarmament & 33/37 & 0.915 & 0.519 & -0.396 & 0.924 & 0.924 & 0.000 \\
Other/institutional & 20/48 & 0.724 & 0.127 & -0.597 & 0.920 & 0.944 & 0.024 \\
\bottomrule
\end{tabularx}
\endgroup
\end{table}


The near-saturation of the 2025 focal snap probabilities raises a natural
diagnostic concern: any terminal year might be classified as a snap because the
model does not observe the following year. Table~\ref{tab:snap-stress-tests}
checks this directly with short stress-test refits. The diagnostic fits use all
theory-selected focal actors plus 45 balanced non-focal actors, so they should be
read as model-pathology triage rather than final inference. The actual 2025
subset still produces high terminal snap probabilities. By contrast, replacing
the 2025 network with a copy of the 2024 network produces no terminal-year
snaps, and treating 2024 or 2023 as the terminal year produces very low focal
snap probabilities. The evidence therefore does not support the interpretation
that 2025 saturation is a generic final-year pattern. The data instead support a
broad annual-agenda re-coordination that is specific to the observed 2025
network.

\clearpage
\begin{table}[!htbp]
\centering
\begingroup
\scriptsize
\caption{Stress tests for terminal-year snap saturation in the UN-voting layer.}
\label{tab:snap-stress-tests}
\begin{tabularx}{\textwidth}{@{}>{\raggedright\arraybackslash}Xrrrrrrr@{}}
\toprule
Test & Year & Actors & Focal obs. & Focal mean & Focal $>.50$ & Focal $>.99$ & All mean \\
\midrule
Main full panel, actual 2025 & 2025 & 200 & 34 & 0.982 & 1.000 & 0.912 & 0.776 \\
Diagnostic subset, actual 2025 & 2025 & 80 & 34 & 0.810 & 0.824 & 0.529 & 0.458 \\
Diagnostic subset, 2024 copied as 2025 & 2025 & 80 & 34 & 0.000 & 0.000 & 0.000 & 0.000 \\
Diagnostic subset, terminal 2024 & 2024 & 80 & 34 & 0.001 & 0.000 & 0.000 & 0.038 \\
Diagnostic subset, terminal 2023 & 2023 & 80 & 34 & 0.029 & 0.029 & 0.029 & 0.059 \\
\bottomrule
\end{tabularx}
\endgroup
\end{table}

\FloatBarrier

The terminal-year checks ask whether a snap model mechanically over-classifies
the last observed year. A separate rolling prediction check asks whether 2025 is
also unusually hard to forecast for a smooth no-snap model. We fit the same
Gaussian AR(1) latent-space model twice: once through 2023 to predict 2024, and
once through 2024 to predict 2025. Table~\ref{tab:oos-prediction} reports
per-state forecast surprise, standardized against each state's own in-sample
one-step forecast history. The contrast is sharp. Predicting 2024 gives a focal
median surprise of $-0.44$, with no focal actors above the high-surprise
threshold. Predicting 2025 gives a focal median surprise of $4.88$, with $94$
percent of observed focal actors above that threshold. The focal 2025 forecast
is also the worst focal one-step forecast in the training comparison set. The
U.S. itself is not the high-surprise unit in this diagnostic; the predictive
break is broad across allies, rivals, and hedgers.

\begin{table}[!htbp]
\centering
\caption{Rolling one-year-ahead prediction diagnostics for the UN-voting layer. Each row fits a smooth AR(1) latent-space model through the training end year and forecasts the next observed annual network. Surprise is the per-state loss in predictive log density relative to that state's in-sample one-step forecast history, scaled by the in-sample median absolute deviation. Higher values mean worse prediction.}
\label{tab:oos-prediction}
\scriptsize
\resizebox{\textwidth}{!}{%
\begin{tabular}{@{}l r r r r r r r r r r@{}}
\toprule
\textbf{Train $\rightarrow$ target} & \textbf{States} & \textbf{Focal} & \textbf{All med.} & \textbf{Focal med.} & \textbf{U.S.} & \textbf{Allies} & \textbf{Rivals} & \textbf{Hedgers} & \textbf{Focal $>2$} & \textbf{Focal worse pct.} \\
\midrule
1989--2023 -> 2024 & 191 & 34 & -0.39 & -0.44 & 1.23 & -0.47 & -0.89 & -0.15 & 0\% & 29\% \\
1989--2024 -> 2025 & 191 & 34 & 3.92 & 4.88 & -6.88 & 5.18 & 4.51 & 4.61 & 94\% & 100\% \\
\bottomrule
\end{tabular}
}
\end{table}


\begin{figure}[ht]
\centering
\includegraphics[width=0.9\textwidth]{oos_prediction_surprise.png}
\caption{Rolling one-year-ahead forecast surprise in the UN-voting layer. Each
column compares a smooth AR(1) no-snap model's forecast for the next annual
network against that state's in-sample one-step forecast history. Higher values
mean worse prediction. Grey points are non-focal states; colored points are the
theory-selected focal actors.}
\label{fig:oos-prediction}
\end{figure}
\FloatBarrier

Figure~\ref{fig:snap-transition} shows another diagnostic: how the posterior
snap probability relates to the size of a transition relative to the fitted AR
drift process. Each grey point is an eligible state-year transition in the
UN-voting panel. The horizontal line marks posterior probability $0.5$; the
vertical line marks one posterior drift standard deviation from the AR prediction.
The relationship is strongly increasing, as it should be, but it is not a
separate descriptive plot of raw movement. The vertical coordinate is the
posterior regime probability from the drift-versus-snap comparison, while the
horizontal coordinate is a posterior-mean diagnostic of transition size. The
model supplies the probability that a transition is better explained as a reset
than as an unusually large continuous innovation.

\begin{figure}[ht]
\centering
\includegraphics[width=0.9\textwidth]{snap_transition_diagnostic.png}
\caption{Transition size and posterior snap probability in the UN-voting fit.
Grey points are all eligible actor-year transitions from 1990 to 2025. The
x-axis reports each transition's distance from its AR-drift prediction in units
of the posterior mean drift standard deviation on a log scale; the y-axis reports
the posterior snap probability. Colored points are the 2025 focal actors, with
selected states labeled.}
\label{fig:snap-transition}
\end{figure}

\subsection{Simulation study}

To verify that the estimand behaves as its interpretation requires, we simulate
networks from the same drift-or-snap transition used in the model, with known
indicators, and check recovery. We generate $n=30$ actors over $T=8$ years with
$R=2$, drawing planted snap years from the diffuse snap component and letting all
other actor-years drift from the previous position. One planted snap is placed in
the terminal year to check the same boundary condition that applies to 2025 in
the empirical application. Figure~\ref{fig:sim-recovery} shows the result:
estimated posterior snap probabilities are low across most non-planted
actor-years and elevated at most planted snaps. Across all planted and
non-planted eligible actor-years, the rank ordering of posterior snap
probabilities yields an area under the ROC curve of approximately $1.00$. The
exercise is a targeted validation check rather than a full Monte Carlo
calibration; the
individual planted-snap probabilities in Table~\ref{tab:validation} show that
some lower-signal planted snaps fall below the $0.5$ interpretive threshold.
Table~\ref{tab:validation} collects the validation quantities.

\begin{table}[ht]
\centering
\caption{Validation of the snap estimand on simulated data. Power is summarized
by the rank separation between planted and non-planted actor-years and by the
individual planted-snap probabilities, including misses. Specificity is checked
on drift-only data. Nesting: with the snap component disabled, the model
reproduces the standard AR(1) dynamic latent-space fit.}
\label{tab:validation}
\small
\begin{tabularx}{0.9\textwidth}{@{}X l@{}}
\toprule
\textbf{Quantity} & \textbf{Value} \\
\midrule
Snap-detection AUC (planted vs.\ non-planted) & $1.00$ \\
Detection probability at the six planted snaps & $0.37,\ 1.00,\ 0.96,\ 1.00,\ 0.29,\ 1.00$ \\
Detection probability at the terminal planted snap & $1.00$ \\
Max.\ false snap probability (recovery run) & $0.090$ \\
\midrule
Max.\ snap probability, drift-only data & $0.263$ \\
Mean snap probability, drift-only data & $0.0057$ \\
Max.\ terminal-year snap probability, drift-only data & $0.133$ \\
Mean terminal-year snap probability, drift-only data & $0.0118$ \\
Posterior snap rate $\pi$, drift-only data & $0.0018$ \\
\midrule
Snaps off ($\pi\to 0$) reproduces AR(1) dynamic latent space & yes \\
\bottomrule
\end{tabularx}
\end{table}

\begin{figure}[ht]
\centering
\includegraphics[width=0.85\textwidth]{sim_recovery.png}
\caption{Recovery on simulated data. Each cell is the posterior snap probability
for a simulated state (rows) in a simulated year (columns); darker is higher.
Planted snaps are marked $\times$. The heatmap reports both recovered planted
snaps and any planted snaps the model fails to separate sharply.}
\label{fig:sim-recovery}
\end{figure}

The complementary property is specificity. Figure~\ref{fig:sim-specificity}
contrasts the distribution of posterior snap probabilities under
\emph{drift-only} data, simulated with no snaps at all, against the
probabilities the model assigns at genuinely planted snaps. Under drift-only
data the probabilities pile up at zero, with a mean of $0.0057$ across eligible
actor-years. The largest drift-only value is $0.263$, so the conservative prior
does not eliminate every isolated false positive; it does, however, make the
drift-only distribution sparse and low. At planted snaps the probabilities are
much higher on average, but not uniformly so, which is why
Table~\ref{tab:validation} reports the individual planted cases. This separation
supports the substantive reading of the results: a near-zero snap
probability is evidence of continuity under this design, while borderline events
should be interpreted with appropriate uncertainty.

\begin{figure}[ht]
\centering
\includegraphics[width=0.85\textwidth]{sim_specificity.png}
\caption{Specificity and power. Top: posterior snap probabilities under
drift-only simulated data (no snaps), which concentrate at zero. Bottom: posterior
snap probabilities at planted snaps, which are generally higher but include weak
recoveries.}
\label{fig:sim-specificity}
\end{figure}

Finally, setting the snap component aside ($\pi\to 0$) reproduces the standard
AR(1) dynamic latent-space fit, confirming that the snap model nests it. Together
these checks support the interpretation of high posterior snap probability as
evidence of discontinuity and near-zero probability as evidence of continuity
under the fitted model.

\section{Robustness Checks}
\label{app:robustness}

We report robustness checks for the main UN-voting result. The baseline is the
$R=2$ snap-shift model reported in the article. We then vary the voting-score
construction, the latent rank, the snap-prior scale, the snap-rate prior, and the
state-composition rule used to summarize or refit the posterior.

\subsection{Voting-score construction}

The main analysis weights UN voting agreement by resolution divisiveness so that
agreement on votes that split the assembly carries more information than
agreement on near-unanimous votes. To check that the 2025 rupture is not an
artifact of this transformation, we refit the snap model on three alternative
UN-voting constructions: unweighted agreement; divisiveness-weighted agreement
after dropping the most divisive ten percent of resolutions within each year; and
unweighted agreement after the same drop.

The 2025 diplomatic rupture is not sensitive to these choices. Under unweighted
agreement, the 2025 focal rupture index is $0.991$ and 2025 ranks first; the
United States, core allies, and pivotal hedgers all have category means of
$1.00$, while strategic rivals average $0.967$. Dropping the most divisive
resolutions makes the result still stronger: the weighted drop-divisive index is
$0.998$, and the unweighted drop-divisive index is $0.996$, both ranked first.
These checks show that the broad 2025 diplomatic rupture is not manufactured by
the divisiveness weighting or by a small set of the most polarizing resolutions.

\subsection{Balanced-state restriction}

Changing state composition could, in principle, make a late-period year appear
unusual because different actors are present. As a check, we recompute the
reported rupture summaries on the $132$ states observed in every UN voting year,
holding the posterior snap probabilities from the main fit fixed. This is a
post-estimation restriction of the reported summaries, not a separate model
refit.

The balanced-state restriction leaves the conclusion unchanged. In the full
summary the focal rupture index rises from $0.030$ in 2024 to $0.956$ in 2025,
making 2025 the highest-ranked year in the series. On the balanced-state
restriction it rises from $0.025$ to $0.979$, and 2025 remains ranked first. The
United States' 2025 snap probability is unchanged at $1.00$ because the United
States is in the balanced panel; the allied-core mean remains $1.00$, the rival
mean is $1.00$, and the hedger mean is $0.939$. The distance pattern also
survives the restriction: allied cohesion remains stable from 2024 to 2025
($0.069$ to $0.068$), while the mean distance from the United States to the allied
core expands from $0.103$ to $0.906$.

\subsection{Historical latent-space comparisons}

The rupture-index calibration is easier to interpret when the latent spaces
around earlier benchmark moments are shown directly. Figure~\ref{fig:hist-space}
compares four before-and-after windows: the post-Cold War transition
($1989$ to $1992$), the post-9/11 and Bush period ($2000$ to $2002$), the first
Trump administration ($2016$ to $2020$), and the 2025 rupture ($2024$ to $2025$).
For each window, the later year is aligned to the earlier year by Procrustes
rotation and translation using states observed in both years. Arrows trace
theory-selected focal actors observed in both years; light grey points show other
observed states; labels mark selected anchor states. In the post-Cold War panel,
the Soviet Union (SUN) and Russia (RUS) are shown as historical anchors rather
than connected as one actor.

The historical comparisons clarify what is distinctive about 2025. Earlier
periods contain movement, but the 2025 panel shows the specific configuration
emphasized in the main text: the United States separates sharply
from the allied cluster, while the allied cluster remains compact. The plot is
descriptive rather than a substitute for the rupture index, because each panel is
scaled to its own before-and-after comparison. Table~\ref{tab:historical-cal}
puts the benchmark moments on the same annual-index scale. For each benchmark we
report the eligible transition-year window, the window mean, the maximum annual
index in that window, the year and post-Cold War rank of that maximum, and the
window maximum as a fraction of the 2025 index.

\begin{table}[ht]
\centering
\caption{Historical calibration of the focal-actor rupture index in UN voting.
All rows use the same annual index and the same window summaries. Ranks are among
eligible transition years, 1990--2025, with 1 the highest focal rupture index.}
\label{tab:historical-cal}
\small
\begin{tabularx}{\textwidth}{@{}X l r r r r r@{}}
\toprule
\textbf{Benchmark} & \textbf{Years} & \textbf{Mean} & \textbf{Max.} &
\textbf{Max. year} & \textbf{Max. rank} & \textbf{Max./2025} \\
\midrule
Post-Cold War transition & 1990--1992 & 0.350 & 0.406 & 1991 & 16 & 0.425 \\
9/11, Afghanistan, and Iraq & 2001--2003 & 0.359 & 0.465 & 2001 & 9 & 0.486 \\
Global financial crisis & 2008--2009 & 0.381 & 0.619 & 2009 & 2 & 0.648 \\
Crimea annexation & 2014--2015 & 0.316 & 0.429 & 2015 & 11 & 0.449 \\
First Trump administration & 2017--2020 & 0.324 & 0.498 & 2019 & 8 & 0.521 \\
Russia full-scale invasion & 2022--2023 & 0.463 & 0.544 & 2023 & 4 & 0.569 \\
Second Trump administration & 2025 & 0.956 & 0.956 & 2025 & 1 & 1.000 \\
\bottomrule
\end{tabularx}
\end{table}

\begin{figure}[ht]
\centering
\includegraphics[width=\textwidth]{historical_space_shifts.png}
\caption{Historical latent-space shifts around benchmark moments in UN voting.
Each panel aligns the later year to the earlier year within the same event
window. Arrows trace focal actors observed in both years; open points mark the
earlier year and filled points mark the later year. Light grey points show other
observed states. Panel scales are free, so distances and arrow lengths should be
read within panels rather than compared visually across panels.}
\label{fig:hist-space}
\end{figure}

\subsection{Transregional allied distances}

The main text argues that the United States separates from the allied bloc
without causing the European and Indo-Pacific wings of that bloc to separate from
one another. Figure~\ref{fig:bridge-distance} tracks this claim over the full
post-Cold War series. The European wing is France, Germany, the United Kingdom,
Italy, the Netherlands, Spain, Belgium, Norway, Denmark, and Poland; Germany is
represented by the observed German state code in each year. The Indo-Pacific
wing is Japan, South Korea, and Australia. Distances are mean pairwise latent
distances within each year, so they are invariant to rotation and translation of
the latent space.

The 2025 change is concentrated in the United States' relation to both allied
wings. The mean distance from the United States to European allies increases from
$0.090$ in 2024 to $0.894$ in 2025. The mean distance from the United States to
Japan, South Korea, and Australia increases from $0.160$ to $0.893$. By contrast,
the mean distance between the European allies and Japan, South Korea, and
Australia is essentially unchanged, $0.107$ to $0.108$. The individual
Indo-Pacific comparisons tell the same story: Japan's mean distance to Europe
falls from $0.171$ to $0.164$, South Korea's from $0.093$ to $0.092$, and
Australia's rises only from $0.058$ to $0.067$. The United States therefore
loses its close position relative to both allied wings, but the distance evidence
does not show a split between Europe and the Indo-Pacific allies.

Figure~\ref{fig:bridge-space} shows the same result as a three-node distance
schematic for 2024 and 2025. The node positions are fixed for readability; the
edge labels are the latent distances. In 2024 the United States is close to both
the European wing and the Japan--South Korea--Australia wing. In 2025 it moves
away from both, while the European and Indo-Pacific wings remain close to one
another.

\begin{figure}[ht]
\centering
\includegraphics[width=0.9\textwidth]{allied_bridge_distances.png}
\caption{Transregional allied distances in the UN-voting latent space.
Top: mean distances from the United States to the European allied wing, from the
United States to Japan, South Korea, and Australia, and between the European and
Indo-Pacific allied wings. Bottom: mean distances from Japan, South Korea, and
Australia to Europe. The dashed line marks 2025.}
\label{fig:bridge-distance}
\end{figure}

\begin{figure}[ht]
\centering
\includegraphics[width=0.9\textwidth]{allied_bridge_space.png}
\caption{Three-node schematic of transregional allied distances in
the UN-voting latent space. Node positions are fixed for readability; edge labels
report the mean latent distances among the United States, European allied wing,
and Japan--South Korea--Australia allied wing. Thicker edges indicate larger
distances.}
\label{fig:bridge-space}
\end{figure}

\begin{table}[ht]
\centering
\caption{Transregional allied distances before and after the 2025 break.
Entries are mean pairwise latent distances in the UN-voting space.}
\label{tab:bridge-distance}
\small
\begin{tabular}{@{}l r r r@{}}
\toprule
\textbf{Distance} & \textbf{2024} & \textbf{2025} & \textbf{Ratio} \\
\midrule
United States to European allies & $0.090$ & $0.894$ & $9.9$ \\
United States to Japan, South Korea, and Australia & $0.160$ & $0.893$ & $5.6$ \\
European allies to Japan, South Korea, and Australia & $0.107$ & $0.108$ & $1.0$ \\
Japan to European allies & $0.171$ & $0.164$ & $1.0$ \\
South Korea to European allies & $0.093$ & $0.092$ & $1.0$ \\
Australia to European allies & $0.058$ & $0.067$ & $1.2$ \\
\bottomrule
\end{tabular}
\end{table}

\subsection{Latent-rank, prior, and panel refits}

The main text uses a conservative rare-snap prior,
$\pi\sim\mathrm{Beta}(1,999)$, and a diffuse snap scale $\kappa=2$. The posterior
snap rate in UN voting is nevertheless high ($0.369$), which should be read as
evidence that the diplomatic layer contains many annual re-coordinations rather
than as a failure of the prior. The point of Table~\ref{tab:un-spec} is narrow:
it asks whether the 2025 conclusion depends on the exact modeling choices used in
the main analysis. If the result disappeared when the latent rank, snap scale,
snap-rate prior, or actor-composition rule changed, then the empirical claim
would be fragile. These refits are sensitivity checks, not replacements for the
main analysis.

\begin{table}[ht]
\centering
\caption{UN-voting sensitivity checks for the 2025 diplomatic rupture. Rows vary
the latent rank, snap scale, snap-rate prior, and actor-composition rule. Entries
are posterior mean snap probabilities in 2025 except for $\bar{\pi}$ and the
index rank. The focal index is the mean across the theory-selected hegemon, core
allies, strategic rivals, and pivotal hedgers.}
\label{tab:un-spec}
\scriptsize
\begin{tabularx}{\textwidth}{@{}l X r r r r r r r@{}}
\toprule
\textbf{Specification} & \textbf{Change} & $\bar{\pi}$ & \textbf{US} &
\textbf{Allies} & \textbf{Rivals} & \textbf{Hedgers} &
\textbf{Index} & \textbf{Rank} \\
\midrule
Main & main analysis fit, $R=2$, $\kappa=2$ & 0.369 & 1.000 & 1.000 & 0.896 & 0.944 & 0.956 & 1 \\
Balanced summary & same fit, $132$ always-observed states & 0.369 & 1.000 & 1.000 & 1.000 & 0.939 & 0.979 & 1 \\
Rank 3 & short refit, $R=3$ & 0.470 & 1.000 & 1.000 & 0.963 & 0.991 & 0.988 & 1 \\
$\kappa=1.5$ & short refit, tighter snap prior & 0.369 & 1.000 & 1.000 & 0.894 & 0.952 & 0.958 & 1 \\
$\kappa=3$ & short refit, wider snap prior & 0.369 & 1.000 & 1.000 & 0.894 & 0.952 & 0.958 & 1 \\
$\pi\sim\mathrm{Beta}(1,1999)$ & short refit, rarer snap prior & 0.314 & 1.000 & 1.000 & 0.894 & 0.947 & 0.956 & 1 \\
Balanced panel & short refit, balanced actor panel & 0.325 & 1.000 & 0.979 & 1.000 & 0.959 & 0.977 & 1 \\
\bottomrule
\end{tabularx}
\end{table}

Across all six checks, 2025 remains the highest-ranked year in the post--Cold War
series. The table rules out a narrow specification explanation: the 2025 result
is not created by the chosen latent dimension, the snap-prior scale, the
rare-snap prior, or late-period changes in actor composition. The rare-prior
refit is especially informative because it doubles the prior odds against snaps
and still leaves the 2025 focal rupture index at $0.956$, with 2025 ranked first.
The balanced-panel refit reaches the same conclusion while estimating the model
only on states observed in every UN voting year.

\section{Alternative Transition Models}
\label{app:alternative-transitions}

The snap-shift model is not defended primarily by global predictive fit. Its
purpose is the actor-year discontinuity estimand. Still, we fit two alternative
transition models on the main UN-voting layer: a Gaussian AR(1) dynamic latent
space with no snap regime, and a Student-$t$ transition that allows large
continuous innovations. Table~\ref{tab:transition-waic} reports WAIC summaries.
The snap model has the best WAIC, but global panel fit is a blunt discriminator:
most dyad-years are structurally stable, so all three models agree on the
overwhelming majority of observations.

\begin{table}[ht]
\centering
\caption{UN-voting transition-model comparison. Higher ELPD and lower WAIC are
better; ELPD differences and standard errors are relative to the snap-shift
model. The comparison is reported as a diagnostic, not as the basis of the
rupture claim.}
\label{tab:transition-waic}
\small
\begin{tabular}{@{}l r r r r r@{}}
\toprule
\textbf{Transition} & \textbf{ELPD WAIC} & \textbf{ELPD diff.} &
\textbf{SE diff.} & \textbf{Eff. parameters} & \textbf{WAIC} \\
\midrule
Snap-shift & 1,250,152 & 0.0 & 0.0 & 949 & -2,500,304 \\
Gaussian AR(1) & 1,249,429 & -722.6 & 27.9 & 860 & -2,498,859 \\
Student-$t$ & 1,249,384 & -768.3 & 25.0 & 873 & -2,498,767 \\
\bottomrule
\end{tabular}
\end{table}

The Student-$t$ fit is the most serious continuous alternative because it can
absorb large movements without assigning a discrete snap. To compare it with the
snap result, we recover each actor-year's AR(1) innovation magnitude from the
fitted Student-$t$ latent paths and compute its percentile among all eligible
actor-years. Table~\ref{tab:t-strain} summarizes the 2025 focal actors by
category. The continuous model also regards 2025 as high strain: the United
States is at the $96.9$th percentile, core allies have a median percentile of
$97.0$, and pivotal hedgers have a median percentile of $91.6$. This supports the
claim that 2025 is unusual even when large continuous innovations are allowed, but
it does not replace the snap estimand because it does not return
$\Pr(\delta_{i,t}=1\mid Y)$.

\begin{table}[ht]
\centering
\caption{Student-$t$ transition strain in 2025. Percentiles compare each focal
actor's 2025 AR(1) innovation magnitude with all eligible actor-years in the
Student-$t$ fit.}
\label{tab:t-strain}
\small
\begin{tabular}{@{}l r r r r@{}}
\toprule
\textbf{Category} & \textbf{N} & \textbf{Mean percentile} &
\textbf{Median percentile} & \textbf{Maximum percentile} \\
\midrule
Hegemon & 1 & 0.969 & 0.969 & 0.969 \\
Core allies & 14 & 0.964 & 0.970 & 0.984 \\
Strategic rivals & 9 & 0.772 & 0.863 & 0.923 \\
Pivotal hedgers & 11 & 0.913 & 0.916 & 0.961 \\
\bottomrule
\end{tabular}
\end{table}

\section{Correlates of the 2025 Snap}
\label{app:whosnaps}

The configuration result is descriptive: it characterizes which states moved and
how the space was reconfigured. As a consistency check, we ask whether pre-2025
covariates help explain the remaining variation in 2025 snap probabilities. This
exercise should be read cautiously for the opposite reason from the old rare-event
framing: the 2025 diplomatic rupture is broad, so the dependent variable is
saturated near one for many states. We report the exercise as descriptive
triangulation, not as a causal account of why any state moved.

Across the $167$ states with complete covariates, $136$ have posterior 2025 snap
probability above $0.5$. States with a United States defense pact have a mean snap
probability of $0.770$, compared with $0.808$ among states without one; formal
alliance status is therefore not the margin that separates high from low snap
probability in this broad event. Table~\ref{tab:who-snaps} reports a
fractional-logit regression of posterior 2025 snap probability on pre-2025 United
States and rival voting alignment, United States and Chinese trade dependence,
trade centrality, prior sanctions exposure, regime type (V-Dem electoral
democracy index), and United States alliance. Continuous predictors are centered
and divided by two standard deviations, so their coefficients are comparable to
the binary alliance indicator. The conventional standard errors and $p$-values
are descriptive diagnostics; they do not propagate uncertainty from the first
stage snap posterior.

\begin{table}[ht]
\centering
\caption{Descriptive fractional-logit correlates of posterior 2025 snap
probability. The dependent variable is the posterior snap probability from the
main UN-voting fit, bounded away from zero and one for estimation. Continuous
predictors are centered and divided by two standard deviations.}
\label{tab:who-snaps}
\small
\begin{tabular}{@{}l r r r@{}}
\toprule
\textbf{Predictor} & \textbf{Estimate} & \textbf{SE} & \textbf{$p$} \\
\midrule
Intercept & $1.734$ & $0.306$ & $<0.001$ \\
United States voting alignment, 2024 & $5.170$ & $2.171$ & $0.018$ \\
Rival voting alignment, 2024 & $3.895$ & $1.973$ & $0.050$ \\
United States trade dependence, 2024 & $-0.762$ & $0.445$ & $0.089$ \\
China trade dependence, 2024 & $0.112$ & $0.487$ & $0.818$ \\
Trade centrality, 2024 & $0.787$ & $0.559$ & $0.161$ \\
Prior sanctions exposure & $0.297$ & $0.584$ & $0.612$ \\
Polyarchy, 2024 & $-1.599$ & $0.599$ & $0.008$ \\
United States defense pact & $0.037$ & $0.625$ & $0.952$ \\
\bottomrule
\end{tabular}
\end{table}

The clearest conditional patterns are a positive association with pre-2025
United States voting alignment and a negative association with polyarchy; rival
alignment is positive but borderline by conventional thresholds. Trade
dependence, trade centrality, prior sanctions exposure, and formal alliance
status are small relative to their uncertainty. The exercise reinforces the
descriptive point: 2025 is too broad a diplomatic rupture for a simple
ally-versus-rival correlate story.


\end{document}
